\section{تدقيق الاصطلاحات وحدود الادعاء}

\subsection{اصطلاح المتغيرات في مسألة وارينغ}

نعتمد في هذا الفصل الاصطلاح الصريح
\[
\mathbb N_{+}=\{1,2,3,\ldots\}.
\]
وعليه فإن عدد التمثيلات $r_{s,k}(N)$ يعد الحلول المرتبة
\[
N=x_1^k+\cdots+x_s^k,
\qquad
x_j\in\mathbb N_{+}.
\]
ويتفق هذا مع مجموع فايل المستعمل في المتن
\[
f_k(\alpha;P)=\sum_{1\le x\le P}e(\alpha x^k).
\]
ولا يدخل الصفر في هذا التعريف. وإذا استعملت نسخة تسمح بـ$x_j=0$ فإن دالة التوليد تتغير بإضافة الحد $1$، ويجب إعادة ضبط عدد التمثيلات والحدود الدنيا، وإن بقي الحد الرئيس الكلاسيكي من الرتبة نفسها في نطاق المتغيرات الكبير.

وللدقة المنتهية نقرأ الشرط $1\le x\le P$ على أنه
\[
1\le x\le \lfloor P\rfloor,
\]
حيث $P=N^{1/k}$. لا يؤثر إدخال الجزء الصحيح في رتبة الحد الرئيس، لكنه يزيل الغموض من الهوية الدقيقة لعد التمثيلات.

\subsection{حدود النتيجة الداخلية}

النتائج المثبتة داخل الفصل هي:
\begin{enumerate}
\item تعامد الدوال الأسية.
\item هوية تكامل عدد التمثيلات.
\item هوية قياس التكامل المفرد بعد التطبيع.
\end{enumerate}
أما تقدير الأقواس الصغرى، والتقريب الكمي الموحد على الأقواس الكبرى، وتقارب وإيجابية السلسلة المفردة في النطاق الكلاسيكي، فتبقى مدخلات مقتبسة. لذلك لا تنسب الصيغة التقاربية الكاملة لمسألة وارينغ إلى الفصل بوصفها مبرهنة مثبتة هنا.

\subsection{حكم التدقيق الاصطلاحي}

\begin{center}
\texttt{TERMINOLOGY-AUDIT-04 = PASS}
\end{center}

أُغلق الغموض المتعلق بـ$\mathbb N$، وثُبت أن مسألة وارينغ في هذا الفصل تستعمل القوى الموجبة فقط، مع قراءة الحد العلوي بواسطة $\lfloor P\rfloor$ في الهويات الدقيقة.
